\documentclass[11pt]{article}
\usepackage[a4paper,margin=25mm]{geometry}
\usepackage{amsmath,amssymb,amsthm,booktabs}
\usepackage[hidelinks]{hyperref}
\emergencystretch=2em
\newtheorem{theorem}{Theorem}
\newcommand{\Z}{\mathbb Z}
\title{A cyclic difference set contradicts\\multiplier-order divisibility\\
\large Disproof of Conjecture 00000006338}
\author{gaochengzhecpu}\date{}
\begin{document}\maketitle
\begin{abstract}
The subset $D=\{3,6,7,12,14\}$ of $\Z/21\Z$ is a nontrivial
$(21,5,1)$ difference set. Its full multiplier group is
$\{1,2,4,8,16,11\}$, of order $6$. This order divides none of
$v=21$, $k=5$, or $\lambda=1$, nor the derived order $k-\lambda=4$
or the product $vk\lambda=105$. Thus the asserted divisibility of the
multiplier-group order by the difference-set parameters fails.
All finite counts and the full multiplier classification are checked in Lean 4.
\end{abstract}

\section*{The necessary assertion being tested}
Conjecture 00000006338 contains the assertion that the order of a
difference-set multiplier group divides the difference-set parameters.
We use their standard meaning $(v,k,\lambda)$: a $k$-element subset
of a group of order $v$ represents each nonzero element exactly
$\lambda$ times as an ordered difference.
It suffices to refute this necessary conjunct; the other assertions
about number-field units and fixed-point counts need not be resolved.

For a cyclic group written additively, the standard multiplier group is
\[
 M(D)=\{m\in(\Z/v\Z)^\times:
          mD=D+a\text{ for some }a\in\Z/v\Z\}.
\]
Thus a multiplier is allowed to send the difference set to a translate.
It is not restricted in advance to fixing the set. All additive
automorphisms of a cyclic group have this numerical form: an additive
map is determined by the image of $1$, and it is an automorphism exactly
when that image is a unit. This agrees with the abelian convention in
\'{O} Cath\'{a}in, Definition 16 and the discussion following it.

\section*{A genuine difference set and its design}
All arithmetic below is modulo $21$. Set
\[
       D=\{3,6,7,12,14\}.
\]
The off-diagonal entries of the following table are the ordered
differences $d-e$:
\[
\begin{array}{c|rrrrr}
 d\backslash e&3&6&7&12&14\\\hline
 3 & -&18&17&12&10\\
 6 & 3&-&20&15&13\\
 7 & 4&1&-&16&14\\
12 & 9&6&5&-&19\\
14 &11&8&7&2&-
\end{array}
\]
Every residue from $1$ to $20$ occurs exactly once. Therefore $D$ is
a $(21,5,1)$ difference set. It is nontrivial, since $20>5>1>0$.

Its developed design has points $\Z/21\Z$ and blocks $D+a$.
These $21$ blocks are distinct: their sums are $5a$, because
$\sum_{d\in D}d=42=0$, and multiplication by $5$ is invertible modulo $21$.
Every block has five points. For distinct points $x,y$, the unique
ordered pair $(d,e)\in D^2$ satisfying $d-e=x-y$ supplies the unique
block with translation $a=x-d=y-e$. Consequently this is a symmetric
$2$-$(21,5,1)$ design, not merely a set with the same cardinality.

\section*{Computing the full multiplier group}
Suppose that a unit $m$ is a multiplier, with $mD=D+a$.
Multiplication by $m$ permutes the elements of $D$ onto those of $mD$.
Summing both sides gives
\[
       0=m\sum_{d\in D}d=\sum_{d\in D}(d+a)=5a.
\]
Since $5$ is a unit modulo $21$, necessarily $a=0$.
Thus in this example every multiplier fixes $D$, a conclusion derived
from the usual translate-allowing definition.

The powers of $2$ modulo $21$ are
\[
       1,2,4,8,16,11,1.
\]
They form a six-element subgroup $H$ of the unit group, and
$2D=\{6,12,14,3,7\}=D$. Hence every element of $H$ fixes $D$.
The full unit group is the disjoint union
\[
\begin{split}
 U&=H\sqcup(-H),\\
 H&=\{1,2,4,8,16,11\},\qquad
 -H=\{20,19,17,13,5,10\}.
\end{split}
\]
These are exactly the residues coprime to $21$.
For each $m\in-H$, one has
\[
       mD=-D=\{7,9,14,15,18\}\ne D.
\]
No such $m$ is a multiplier, because any allowed translation was
already proved to be zero. This proves $M(D)=H$ and $|M(D)|=6$.

\begin{theorem}
The multiplier-order divisibility assertion in Conjecture 00000006338
is false.
\end{theorem}
\begin{proof}
For the difference set above, the full multiplier group has order $6$,
whereas
\[
       6\nmid21,\qquad 6\nmid5,\qquad 6\nmid1.
\]
Thus the order divides none of the standard parameters. Even if one
also considers the derived order or the product of parameters,
$6\nmid(5-1)=4$ and $6\nmid21\cdot5\cdot1=105$.
The claimed universal divisibility therefore fails on this nontrivial
cyclic difference-set design.
\end{proof}

\section*{Formal correspondence and reproduction}
The Lean project uses Mathlib's ring \texttt{ZMod 21}, finite sets,
unit group, additive homomorphisms and subgroups. The function
\texttt{differenceCount} counts actual ordered pairs with a prescribed
difference. The development into blocks is explicitly constructed;
Lean checks its block count, block sizes and unique pair incidence.

The predicate \texttt{IsMultiplier} includes coprimality and an
existential translation. The complete classification is checked over
all residues and all translations. The finite set of six units is
given its actual subgroup structure, and
\texttt{group\_is\_full\_multiplier\_group} proves that its membership
is equivalent to the defining multiplier predicate.
\texttt{all\_automorphism\_multipliers} also verifies the correspondence
for arbitrary additive automorphisms, using their value at $1$.
The order is the cardinality of this subgroup, not a number assigned
by a separate definition. The final theorem combines the difference-set
property, full group identification, order and failed divisibilities.

Finite checks use kernel-checked decision procedures. No external
numerical oracle is used. The accompanying \texttt{verify.py} regenerates
the tables and counts as a supplementary check; the Lean proof does
not depend on its output. The author performed a separate adversarial
self-review. No independent reviewer or subagent was used.

The project pins Lean 4.19.0 and Mathlib revision\\
\texttt{c44e0c8ee63ca166450922a373c7409c5d26b00b}.
In \texttt{lean/}, run
\begin{verbatim}
lake exe cache get
lake build
lake env lean -DwarningAsError=true Main.lean
\end{verbatim}
The dependency cache command is needed only when the pinned official
dependencies are absent. The submitted project itself is rebuilt in
a fresh directory. Run \texttt{python verify.py} for the supplementary
calculation and \texttt{tectonic main.tex} for the PDF. Exact source,
build logs, axiom dependencies, hashes and review records accompany
the submission.

\section*{References}
\begin{enumerate}
\item The Last Math Competition, Conjecture 00000006338.
The exact bilingual source is reproduced in \texttt{SOURCE.md}.
\item P. \'{O} Cath\'{a}in, \emph{Inequivalence of difference sets:
On a remark of Baumert}, Electronic Journal of Combinatorics 20(1)
(2013), P38, Definition 16 and its following discussion.
\href{https://www.combinatorics.org/ojs/index.php/eljc/article/download/v20i1p38/pdf/}{Publisher's full text}.
\end{enumerate}
\end{document}
